\documentclass[11pt,a4paper]{article}

% ---------- Packages ----------
\usepackage[margin=1in]{geometry}
\usepackage{amsmath, amssymb, amsthm, mathtools}
\usepackage{microtype}
\usepackage{algorithm}
\usepackage{algpseudocode}
\usepackage[colorlinks=true,
            linkcolor=blue!50!black,
            urlcolor=blue!50!black,
            citecolor=blue!50!black]{hyperref}

% ---------- Theorem environments ----------
\theoremstyle{plain}
\newtheorem{theorem}{Theorem}[section]
\newtheorem{proposition}[theorem]{Proposition}
\newtheorem{lemma}[theorem]{Lemma}
\newtheorem{corollary}[theorem]{Corollary}

\theoremstyle{definition}
\newtheorem{definition}[theorem]{Definition}

\theoremstyle{remark}
\newtheorem{remark}[theorem]{Remark}
\newtheorem{example}[theorem]{Example}

% ---------- Math shortcuts ----------
\newcommand{\R}{\mathbb{R}}
\newcommand{\N}{\mathbb{N}}
\newcommand{\Z}{\mathbb{Z}}
\newcommand{\E}{\mathbb{E}}
\newcommand{\Var}{\operatorname{Var}}
\newcommand{\Cov}{\operatorname{Cov}}
\newcommand{\F}{\mathcal{F}}
\newcommand{\Prob}{\mathbb{P}}
\newcommand{\Q}{\mathbb{Q}}
\newcommand{\dd}{\mathop{}\!\mathrm{d}}
\newcommand{\eqdef}{\overset{\mathrm{def}}{=}}
\newcommand{\indic}{\mathbf{1}}
\DeclareMathOperator{\Disc}{D^*}

% ---------- Document metadata ----------
\title{Quasi-Monte Carlo: Foundations}
\author{Quant Finance Portfolio --- Phase 2, Block 3.0}
\date{\today}

\begin{document}
\maketitle

\begin{abstract}
We introduce the framework of quasi-Monte Carlo (QMC) integration:
deterministic sequences of points designed to fill the unit cube
more uniformly than random samples, achieving asymptotic error
$O((\log n)^d / n)$ in dimension $d$ versus the standard
$O(1/\sqrt n)$ of Monte Carlo. The price is the loss of statistical
machinery (no central limit theorem, no confidence intervals); we
recover this via randomized QMC (RQMC) with digital shifting. We
develop the discrepancy framework, prove the Koksma-Hlawka
inequality, and construct two foundational sequences: Halton (an
extension of van der Corput's one-dimensional sequence) and Sobol
(based on primitive polynomials over $\mathbb{F}_2$). The relation
to inverse sampling --- which we adopted in Block~0 of this phase
specifically for QMC compatibility --- is highlighted throughout.
\end{abstract}

\tableofcontents

\section{From Monte Carlo to quasi-Monte Carlo}
\label{sec:mc-to-qmc}

The Monte Carlo identity is
\begin{equation}
\int_{[0,1]^d} f(u) \dd u \;\approx\; \frac{1}{n} \sum_{i=1}^n f(u_i),
\qquad u_i \sim \mathrm{Unif}\bigl([0,1]^d\bigr) \text{ i.i.d.},
\end{equation}
with statistical error of order $\sigma(f)/\sqrt n$ where
$\sigma^2(f) = \Var(f(u))$. The variance reduction techniques of
Block~2 attack the constant $\sigma(f)$; they leave the rate
$1/\sqrt n$ untouched.

QMC takes a different route. Instead of choosing $\{u_i\}$
randomly, we choose them deterministically to be ``as uniform as
possible'' in $[0,1]^d$. The hope is that the average $\frac{1}{n}
\sum f(u_i)$ converges to the integral faster than $1/\sqrt n$.

The challenge is to make ``as uniform as possible'' precise. The
notion is captured by the \emph{discrepancy} of a point set, which
measures the worst-case deviation of the empirical distribution
from the uniform distribution on rectangles.

\begin{definition}[star discrepancy]
\label{def:discrepancy}
For a finite point set $P = \{u_1, \ldots, u_n\} \subset [0,1)^d$
and a rectangle $B = \prod_{j=1}^d [0, b_j) \subset [0, 1)^d$, let
\begin{equation}
\Disc(P; B) \;=\; \left| \frac{\#\{i : u_i \in B\}}{n} -
                         \mathrm{vol}(B) \right|.
\end{equation}
The \emph{star discrepancy} of $P$ is
\begin{equation}
\Disc(P) \;=\; \sup_{B \subset [0,1)^d} \Disc(P; B).
\end{equation}
\end{definition}

A sequence with small discrepancy distributes points evenly across
all anchored rectangles. For random points, $\E[\Disc(P)] =
\Theta(\sqrt{(\log \log n) / n})$ by the law of the iterated
logarithm; for QMC sequences, deterministic constructions can
achieve $\Disc(P) = O((\log n)^d / n)$.

\begin{theorem}[Koksma-Hlawka inequality]
\label{thm:kh}
For any function $f : [0,1]^d \to \R$ of bounded variation in the
sense of Hardy-Krause, with variation $V(f)$, and for any finite
point set $P \subset [0, 1)^d$,
\begin{equation}
\left| \frac{1}{n} \sum_{i=1}^n f(u_i)
       \,-\, \int_{[0,1]^d} f(u) \dd u \right|
\;\leq\; V(f) \cdot \Disc(P).
\end{equation}
\end{theorem}

The Koksma-Hlawka bound is \emph{deterministic}: no probability is
involved. Its strength is that, for nice $f$ (smooth, bounded
mixed partial derivatives), $V(f) < \infty$, and the integration
error decays as fast as the discrepancy. Its weakness, on which we
return below, is that $V(f)$ is typically very hard to bound (and
often infinite for the indicators that appear in option payoffs),
so the inequality gives no usable a priori error estimate.

\subsection{The asymptotic comparison}

For Halton or Sobol sequences in dimension $d$, the discrepancy
satisfies $\Disc(P_n) = O((\log n)^d / n)$. Compared to MC:
\begin{center}
\begin{tabular}{l c c}
Method & Error rate & Behaviour as $n \to \infty$ \\
\hline
MC                & $O(n^{-1/2})$         & always wins eventually \\
QMC, $d$ moderate & $O((\log n)^d / n)$  & wins MC for moderate $n$ \\
QMC, $d$ large    & $O((\log n)^d / n)$  & loses to MC for practical $n$
\end{tabular}
\end{center}

The crossover point depends on $d$. For $d \leq 5$, QMC dominates
MC at any $n \geq 10^3$. For $d = 30$, the $(\log n)^d$ factor is
catastrophic --- $(\log 10^6)^{30} \approx 10^{34}$ --- and QMC
requires extreme $n$ to outperform MC. In practice, QMC is most
useful for low-to-moderate dimensional problems.

A modern viewpoint, developed by Sloan, Wo\'zniakowski, Owen and
others, refines this picture: many financial integrands, despite
high nominal dimension, have ``effective dimension'' much smaller
than $d$ because the function depends primarily on a low-dimensional
subspace. QMC then performs well even at large nominal $d$. This is
beyond the scope of the present writeup but motivates the
preparatory experiments of Block~3.1.

\section{The van der Corput sequence}
\label{sec:vdc}

The simplest QMC construction is one-dimensional. Given an integer
base $b \geq 2$, write each non-negative integer $i$ in base $b$ as
\begin{equation}
i \;=\; \sum_{k \geq 0} a_k(i) \, b^k,
\qquad a_k(i) \in \{0, 1, \ldots, b-1\}.
\end{equation}
The \emph{radical inverse function} $\phi_b : \N \to [0, 1)$ reverses
the digits and places a decimal point:
\begin{equation}
\phi_b(i) \;=\; \sum_{k \geq 0} a_k(i) \, b^{-(k+1)}.
\end{equation}
The \emph{van der Corput sequence in base $b$} is $\{\phi_b(i)\}_{i
\geq 1}$.

\begin{example}[base 2]
The base-2 van der Corput sequence starts
\[
\tfrac{1}{2}, \;\tfrac{1}{4}, \;\tfrac{3}{4}, \;\tfrac{1}{8},
\;\tfrac{5}{8}, \;\tfrac{3}{8}, \;\tfrac{7}{8}, \;\tfrac{1}{16},
\;\tfrac{9}{16}, \ldots
\]
After $2^k$ points, the sequence has used each ``odd numerator over
$2^{k+1}$'' exactly once, so the points are perfectly equidistributed
in $[0, 1)$ at scale $2^{-(k+1)}$.
\end{example}

The van der Corput sequence in base $b$ has discrepancy
$\Disc = O(\log n / n)$, optimal in $d = 1$.

\section{The Halton sequence}
\label{sec:halton}

The natural multi-dimensional extension is to use a different prime
base in each coordinate. Choose primes $p_1 < p_2 < \cdots < p_d$
(e.g., $p_1 = 2, p_2 = 3, p_3 = 5, \ldots$), and define the Halton
sequence by
\begin{equation}
\mathbf{u}_i \;=\; \bigl(\phi_{p_1}(i), \phi_{p_2}(i), \ldots,
                         \phi_{p_d}(i)\bigr) \in [0, 1)^d.
\end{equation}
The construction is constructive, infinite, and trivial to
implement: for each new index $i$, evaluate $d$ radical inverses.

The Halton sequence has discrepancy
$\Disc = O((\log n)^d / n)$. The constant in the big-O grows with
$d$: more precisely, the asymptotic constant is $\prod_{j=1}^d
\frac{p_j - 1}{2 \log p_j}$, which for $p_d$ around $20$ (i.e., $d
\approx 8$) becomes uncomfortably large.

\paragraph{The high-dimensional pathology.} Halton's reliance on
primes that grow with dimension introduces a well-known artefact:
for $d \geq 7$, projections of the Halton sequence onto pairs of
high-index dimensions exhibit visible linear correlations. The
prime $p_8 = 19$ is large enough that the radical inverse
``advances'' very slowly within a window of consecutive indices,
producing alignment with $p_9 = 23$. This is why Sobol, despite
being more complex to construct, is preferred in practice.

\paragraph{Implementation.} The radical inverse is computed
iteratively: while $i > 0$, extract the lowest digit of $i$ in base
$b$, multiply by $b^{-(k+1)}$, accumulate. The whole Halton
sequence in dimension $d$ requires $O(d \log_2 n)$ digit operations
per point.

\section{The Sobol sequence}
\label{sec:sobol}

Sobol sequences are the standard choice for QMC in finance.
Whereas Halton uses a different prime base per dimension, Sobol
uses base $2$ in every dimension and varies a different ingredient:
the \emph{direction numbers}, which encode primitive polynomials
over $\mathbb{F}_2$.

\subsection{Direction numbers and the recursion}

For each dimension $j$, fix a sequence of \emph{direction numbers}
$v_1^{(j)}, v_2^{(j)}, \ldots$ where each $v_k^{(j)} \in (0, 1)$.
The Sobol value for index $i$ in dimension $j$ is
\begin{equation}
\label{eq:sobol-construction}
u_i^{(j)} \;=\; b_1(i) \, v_1^{(j)} \,\oplus\, b_2(i) \, v_2^{(j)}
\,\oplus\, \cdots \,\oplus\, b_K(i) \, v_K^{(j)},
\end{equation}
where $b_k(i)$ are the binary digits of the index $i$ ($i = \sum_k
b_k(i) 2^{k-1}$), $K$ is the number of bits needed to represent
$i$, and $\oplus$ is bitwise XOR (interpreted as addition modulo $2$
on the binary expansion of $v_k^{(j)}$).

For the construction to produce a low-discrepancy sequence, the
direction numbers must satisfy specific algebraic properties. They
are derived from a primitive polynomial over $\mathbb{F}_2$ for
each dimension; the construction is due to Sobol (1967) and refined
by Joe and Kuo (2008) into a tabulated form.

\subsection{The Joe-Kuo tables}

Joe and Kuo published direction-number tables for dimensions up to
$21\,201$ in their 2008 paper. The tables are public and freely
distributed at the URL
\href{https://web.maths.unsw.edu.au/~fkuo/sobol/}{web.maths.unsw.edu.au/$\sim$fkuo/sobol/}.
They are the standard reference used by virtually all production
implementations of Sobol in finance (including \texttt{scipy.stats.qmc.Sobol},
which loads them at import time).

The structure of the table is, for each dimension $d$:
\begin{itemize}
\item A primitive polynomial degree $s$ and an integer encoding $a$
representing the polynomial coefficients.
\item Initial direction numbers $m_1, m_2, \ldots, m_s$ (one per
degree-of-freedom of the polynomial).
\end{itemize}
The remaining direction numbers are generated by the recursion
\begin{equation}
m_k \;=\; 2 a_1 m_{k-1} \oplus 2^2 a_2 m_{k-2} \oplus \cdots \oplus
          2^{s-1} a_{s-1} m_{k-s+1} \oplus (2^s m_{k-s} \oplus
          m_{k-s}),
\quad k > s,
\end{equation}
where $a_j$ are the polynomial coefficients (with $a_0 = a_s = 1$
by convention). The direction numbers $v_k = m_k / 2^k$ are then
inserted in equation~\eqref{eq:sobol-construction}.

The construction may look intimidating; in practice, it is~$\sim 50$
lines of code once the Joe-Kuo file is loaded.

\subsection{The Gray-code optimisation}

A naive implementation evaluates the XOR over $K$ direction numbers
for each new index $i$. Sobol observed that consecutive indices in
\emph{Gray-code order} differ in exactly one bit, so generating
$u_{i+1}$ from $u_i$ requires only one XOR. This optimisation is
essential for performance; we adopt it in the C++ implementation.

\section{Quasi-Monte Carlo with inverse sampling}
\label{sec:qmc-inversion}

Suppose we want to estimate $\E[g(Z)]$ where $Z \sim \mathcal{N}(0,
1)^d$. The MC route is:
\begin{enumerate}
\item Draw $u_i \sim \mathrm{Unif}([0,1)^d)$ i.i.d.
\item Apply $z_i = \Phi^{-1}(u_i)$ componentwise.
\item Average $g(z_i)$.
\end{enumerate}
The QMC route replaces step 1 by a low-discrepancy sequence and
keeps steps 2-3 unchanged. \emph{Crucially, this only works if the
mapping $\Phi^{-1}$ is monotone and one-to-one}. Other normal
samplers --- Box-Muller (which mixes pairs of uniforms in a
non-monotone way), Marsaglia polar (which involves rejection),
Ziggurat (rejection-based) --- destroy the low-discrepancy structure
of the input. Box-Muller, in particular, can introduce spurious
correlations between dimensions when paired with a QMC source.

This is why we adopted inversion sampling in Block~0 of this phase,
even though Ziggurat is faster for plain MC. The pipeline
\begin{equation}
\text{QMC source} \;\longrightarrow\; \Phi^{-1}
\;\longrightarrow\; \text{normal samples}
\end{equation}
preserves the discrepancy structure across the transformation. The
choice of normal sampler made in Block~0 was deliberately taken
with this future need in mind.

\section{The need for randomization: from QMC to RQMC}
\label{sec:rqmc}

Consider the situation after running a QMC integrator. We have
computed
\begin{equation}
\hat I_n^{\mathrm{QMC}} \;=\; \frac{1}{n} \sum_{i=1}^n f(u_i),
\end{equation}
a single deterministic number. The Koksma-Hlawka bound gives a
deterministic error estimate $|I - \hat I_n^{\mathrm{QMC}}| \leq
V(f) \Disc$, but $V(f)$ is essentially never available: the
Hardy-Krause variation of an option payoff is hard to bound, often
infinite for indicator functions, and not available in closed form
even for the European call.

\emph{We have a number with no error bar}. This is unacceptable in
production: pricing requires a quantifiable uncertainty.

\subsection{Digital shifting}

The fix is to randomise the QMC sequence in a way that preserves
its low-discrepancy property but makes each marginal point uniform.
The simplest randomisation is the \emph{digital shift}.

\begin{definition}[digital shift]
\label{def:digital-shift}
Given a base-$2$ QMC sequence $\{u_i\}$ in $[0, 1)^d$ and a uniform
random shift $\xi \in [0, 1)^d$, the digitally shifted sequence is
\begin{equation}
\tilde u_i^{(j)} \;=\; u_i^{(j)} \oplus \xi^{(j)}
\quad \text{(bitwise XOR)}.
\end{equation}
For each fixed $i$ and $j$, $\tilde u_i^{(j)} \sim
\mathrm{Unif}([0, 1))$ marginally; the dependence structure between
indices is preserved.
\end{definition}

The digital shift is reversible (XOR with $\xi$ twice gives back the
original), preserves the local properties of Sobol points (by
linearity of XOR over $\mathbb{F}_2$), and makes each individual
$\tilde u_i$ marginally uniform.

\subsection{The RQMC estimator}

\begin{definition}[RQMC estimator with $R$ replications]
\label{def:rqmc}
Choose $R \geq 2$ replications. For each replication $r = 1, \ldots,
R$, draw an independent random shift $\xi^{(r)}$ and compute
\begin{equation}
\hat I_n^{(r)} \;=\; \frac{1}{n} \sum_{i=1}^n f(u_i \oplus \xi^{(r)}).
\end{equation}
The RQMC estimator is the average across replications,
\begin{equation}
\hat I_{n, R}^{\mathrm{RQMC}} \;=\; \frac{1}{R} \sum_{r=1}^R \hat I_n^{(r)}.
\end{equation}
\end{definition}

The $R$ replication estimates $\{\hat I_n^{(r)}\}_{r=1}^R$ are
i.i.d., so we have a CLT for them. The RQMC estimator's variance is
$\sigma_{\mathrm{rep}}^2 / R$ where $\sigma_{\mathrm{rep}}^2 =
\Var(\hat I_n^{(r)})$, and the half-width is
\begin{equation}
h_{n, R}^{\mathrm{RQMC}} \;=\; z_{\beta/2} \cdot
\frac{\hat\sigma_{\mathrm{rep}}}{\sqrt R},
\end{equation}
identical in form to the MC half-width but with the variance estimated
from $R$ replication-estimates instead of from $n$ payoffs.

The crucial empirical fact: $\sigma_{\mathrm{rep}}^2$ scales with
the QMC error, so for nice integrands,
\begin{equation}
\sigma_{\mathrm{rep}}^2 \;=\; O\bigl((\log n)^{2d} / n^2\bigr)
\quad\text{instead of}\quad O(\sigma^2(f) / n).
\end{equation}
At fixed total budget $nR$, RQMC produces estimator variance
asymptotically smaller than MC by orders of magnitude (for nice $f$
and moderate $d$).

\subsection{The trade-off in choosing $n$ vs $R$}

For a fixed total budget $N = nR$ payoff evaluations:
\begin{itemize}
\item Larger $n$ exploits the QMC convergence rate.
\item Larger $R$ gives more replications and a tighter half-width
estimate (improved by $1/\sqrt R$ in the CLT sense, but also useful
for stability of $\hat\sigma_{\mathrm{rep}}$).
\end{itemize}
The standard practical rule is $R \in [10, 30]$ replications and the
remaining budget on $n$. We adopt $R = 20$ in the validation script
of Block~3.1.

\section{Summary and look ahead to Block 3.1}
\label{sec:summary}

We have established the framework:
\begin{itemize}
\item QMC sequences (Halton, Sobol) achieve discrepancy $O((\log n)^d
/ n)$, faster than the MC rate $O(\sqrt{(\log\log n)/n})$.
\item The Koksma-Hlawka bound is deterministic but practically
unusable for option payoffs because $V(f)$ is hard to bound.
\item Inverse sampling (chosen in Block~0) is the only normal
sampler that preserves QMC structure.
\item RQMC with digital shifting recovers a usable error bar by
averaging $R$ independent shifts.
\end{itemize}

Block~3.1 specialises this framework to the European call. We will
implement Halton and Sobol generators, build a multi-step Euler MC
pricer with $d = 20$ (one normal per discretisation step), compare
deterministic-QMC and RQMC against the IID Block~1.2.1 baseline,
and measure empirical convergence rates on a log-log slope plot.

\section{Bibliographic notes}
\label{sec:biblio}

The standard reference for QMC integration is Niederreiter's monograph
\cite{Niederreiter}. The original Sobol sequence is from
\cite{Sobol1967}. The Joe-Kuo tables and primitive polynomials
construction are in \cite{JoeKuo2008}; the table file
\texttt{new-joe-kuo-6.21201} is hosted at \cite{JoeKuoSite}.
Owen's work on RQMC, including the more sophisticated nested
scrambling, is reviewed in \cite{Owen2003}. The financial-applications
viewpoint, including effective dimension, is developed by Glasserman
\cite[Chapter~5]{Glasserman}.

\begin{thebibliography}{99}

\bibitem{Glasserman}
P.\ Glasserman,
\emph{Monte Carlo Methods in Financial Engineering},
Springer, 2004.

\bibitem{JoeKuo2008}
S.\ Joe and F.\ Y.\ Kuo,
\emph{Constructing Sobol' sequences with better two-dimensional
projections},
SIAM Journal on Scientific Computing, 30(5):2635--2654, 2008.

\bibitem{JoeKuoSite}
S.\ Joe and F.\ Y.\ Kuo,
\emph{Sobol' sequence generator},
\url{https://web.maths.unsw.edu.au/~fkuo/sobol/}.

\bibitem{Niederreiter}
H.\ Niederreiter,
\emph{Random Number Generation and Quasi-Monte Carlo Methods},
SIAM, 1992.

\bibitem{Owen2003}
A.\ B.\ Owen,
\emph{Quasi-Monte Carlo sampling},
in \emph{Monte Carlo Ray Tracing}, ACM SIGGRAPH course notes, 2003.

\bibitem{Sobol1967}
I.\ M.\ Sobol',
\emph{On the distribution of points in a cube and the approximate
evaluation of integrals},
USSR Computational Mathematics and Mathematical Physics,
7(4):86--112, 1967.

\end{thebibliography}

\end{document}
